\section{Limitations, Falsifiability, and Research Agenda}

\subsection{Identification limits}

Acceleration is difficult to identify from short, revised, selected, and heterogeneous series. A change point can be mistaken for smooth acceleration; a logistic curve can look exponential far below its ceiling; an accelerating curve can mimic bounded growth over a short window. Latent capability, cost, reliability, automation, adoption, and deployment are endogenous and often measured on different task populations. Elasticities in Equation~\eqref{eq:technical-growth} may be weakly identified and correlated. The generalized realization family adds flexibility but also additional identification burden; sensitivity analysis cannot substitute for data.

\subsection{Measurement and comparability}

A recent source can contain an old measurement, and a familiar label can conceal a changed construct. Benchmark suites, scaffolds, elicitation effort, pricing units, business-survey wording, and target populations evolve. The BTOS audit in Section~8 shows that a change from ``producing goods or services'' to ``any business function'' creates a measurement break that should not be interpreted as pure adoption growth. The protocol can flag such breaks, but a valid bridge requires overlapping measurements or a defensible linking model.

\subsection{External validity}

Software tasks with clear tests differ from organizations, scientific laboratories, physical infrastructure, care systems, and regulated decisions. Benchmark task horizon is not wall-clock runtime, and low-context human baselines may not match resident experts. Scaffolds, tools, security constraints, tacit context, and messy environments can materially alter performance. The framework requires domain-specific task graphs and evidence; it is not a universal coefficient table.

\subsection{Distributional uncertainty}

Residual bootstrap intervals inherit the fitted model pool, residual stationarity assumptions, sample size, and historical support. They are not Bayesian posteriors and do not automatically capture unknown unknowns, source revisions, structural breaks outside the candidate set, or strategic behavior. Six METR hindcast origins are far too few for strong calibration claims. Scenario probabilities remain judgmental unless supported by a validated elicitation or event model.

\subsection{Discontinuity uncertainty}

A recursive-improvement regime may not occur, may arrive outside the horizon, may be constrained by evaluation and grounding, or may improve narrow research loops without general transfer. Conversely, sparse monitoring could detect it late. The hazard structure is therefore conditional, not a calibrated probability in the absence of substantial transition data. Double-exponential paths are stress models, not claims of infinite physical output.

\subsection{Prospective and independent validation}

The local package can verify schemas, algorithms, builds, checksums, and frozen demonstrations. It cannot self-certify future predictive accuracy, independent replication, peer review, legal approval, or an external publication timestamp. A live prospective registry becomes scientifically meaningful only when forecasts are committed before outcomes, matured without retroactive editing, and scored by independent parties.

\subsection{Normative and legal limits}

The decision objective embeds values through utility, loss, downside aversion, authority, and constraints. It cannot decide those values neutrally. Forecasts in finance, medicine, law, safety, regulation, employment, or public policy may create professional and statutory duties beyond this research protocol. Repository notices reduce ambiguity but do not eliminate duties, claims, or jurisdiction-specific requirements. Patent, authorship, licensing, privacy, data-use, and publication decisions require qualified review where material.

\subsection{Falsifiable expectations}

The framework should be revised or rejected if, across preregistered targets:
\begin{enumerate}
\item it fails to improve proper scores, calibration, interval coverage, or decision regret over simpler baselines;
\item the canonical branches add complexity without predictive, diagnostic, or decision value;
\item trigger-based updating does not detect breaks faster than fixed schedules;
\item task-graph estimates systematically miss verification, fallback, authority, and external-waiting time;
\item model averaging is persistently worse than robust single-model rules;
\item the generalized family changes decisions mainly through unidentified parameter freedom;
\item independent teams cannot reproduce the controlled and frozen-reference results from the released source;
\item registered forecasts are systematically overconfident or fail to improve after recalibration.
\end{enumerate}

\subsection{Research program}

Priority work includes preregistered rolling hindcasts over multiple targets; reliability-adjusted cost datasets; authority and adoption measurement; bridge studies across changing constructs; causal or bounded estimates of task-level elasticities; hierarchical models across domains; formal double-counting diagnostics; correlated-failure models; transition-hazard elicitation; environment-stable containers; and evaluation of decisions rather than predictions alone. Independent replication, adversarial review, and a public maturity-and-scoring process are essential.
